\documentclass{article}
\usepackage{fuzz}
\makeatletter
\newenvironment{theorem}[1]{\medskip\noindent\textbf{Theorem (#1).}\@nameuse{zed}}{\@nameuse{endzed}}
\makeatother

\begin{document}

\title{A Formal Specification of \texttt{model-provider-claude} (Claude Code only)}
\author{Punt Labs}
\maketitle

\section{Purpose and scope}

This is the other half of a two-tool split. \texttt{model-provider-claude}
configures exactly one harness --- Claude Code --- and carries all the
complexity that, in the combined tool, made the whole thing fragile:
canonicalization, the \texttt{modelOverrides} map, the tier-alias fan-out, the
small-fast background model, the $1M$-context marker, and the one legitimate use
of a Claude subscription credential. None of that leaks into the sibling
\texttt{model-provider-fleet} tool; it is quarantined here, where it can be
stated precisely and checked hard, because it only ever concerns one harness.

The guarantees this tool must make, which the combined tool only ever got right
by accident:
the main request carries the selection; each tier alias (the Agent tool's
\texttt{opus|sonnet|haiku|fable} sub-agents) resolves to a tier-appropriate id
in the selection's provider family --- a haiku sub-agent runs haiku, not a
force-upgrade; the small-fast background model is pinned to a routable id;
Claude Code's canonicalization is corrected by an override map so the selection
reaches the wire verbatim rather than a mis-canonicalized base; and a
non-Claude selection is refused loudly rather than silently mis-routed.

\section{Given sets and enumerations}

A model id is an opaque catalog token (e.g.
\texttt{openrouter/anthropic/claude-opus-4.8}). A provider is the upstream a
gateway routes to. A canonical id is the base name Claude Code derives from a
model id by dropping the provider prefix and rewriting dots to hyphens (e.g.
\texttt{claude-opus-4-8}).

\begin{zed}
  [ModelId, Provider, CanonId]
\end{zed}

Claude Code issues three request kinds. \emph{main} is the primary loop;
\emph{tier} is an Agent-tool sub-agent spawn resolved through a tier alias;
\emph{small} is the background small-fast model.

\begin{zed}
  ReqKind ::= main | tier | small
\end{zed}

The gateways; the spec is written over an abstract gateway, these are the
instances. $direct$ is modelled separately (absence of a gateway).

\begin{zed}
  Gateway ::= bifrost | envoy | agentgateway
\end{zed}

Authentication: the per-host gateway virtual key, or the operator's Claude
subscription OAuth passed through the gateway. Unlike the fleet tool, $subauth$
has a legitimate home here --- this is Claude Code, Anthropic's own client.

\begin{zed}
  AuthMode ::= vkauth | subauth
\end{zed}

Booleans.

\begin{zed}
  Bool ::= yes | no
\end{zed}

\section{The gateway contract}

$gwCatalog$ gives, per id, the provider; $gwServes$ the entitled subset;
$gwCanon$ the canonical id Claude Code derives from a served id; $gwAnthropic$
whether the gateway exposes the Anthropic protocol Claude Code speaks. Every
gateway entitles at least one id and canonicalizes exactly its catalogued ids.

\begin{axdef}
  gwCatalog : Gateway \fun (ModelId \pfun Provider) \\
  gwServes : Gateway \fun \power ModelId \\
  gwCanon : Gateway \fun (ModelId \pfun CanonId) \\
  gwAnthropic : Gateway \fun Bool
\where
  \forall g : Gateway @ \\
  \t1 \dom (gwCanon~g) = \dom (gwCatalog~g) \land \\
  \t1 gwServes~g \subseteq \dom (gwCatalog~g) \land \\
  \t1 gwServes~g \neq \emptyset
\end{axdef}

The selectable ids are the entitled set.

\begin{axdef}
  validOn : Gateway \fun \power ModelId
\where
  \forall g : Gateway @ validOn~g = gwServes~g
\end{axdef}

Provider of an id on a gateway, for catalogued ids.

\begin{axdef}
  provOn : Gateway \cross ModelId \pfun Provider
\where
  \forall g : Gateway; m : ModelId @ \\
  \t1 ( ( g, m ) \in \dom provOn \iff m \in \dom (gwCatalog~g) ) \land \\
  \t1 ( m \in \dom (gwCatalog~g) \implies provOn~(g, m) = gwCatalog~g~m )
\end{axdef}

A gateway can carry Claude Code iff it exposes the Anthropic protocol.

\begin{axdef}
  usableGateways : \power Gateway
\where
  usableGateways = \{ g : Gateway | gwAnthropic~g = yes \}
\end{axdef}

\section{Which ids Claude Code can run}

Claude Code is Claude-only: it accepts only Anthropic-family model ids. $isClaude$
is the predicate deciding this --- a fixed subset of $ModelId$. A selection
outside it is refused at the top level (the loud \texttt{unrecognized\_model}
outcome), never configured.

\begin{axdef}
  claudeIds : \power ModelId
\where
  claudeIds \neq \emptyset
\end{axdef}

\section{Canonicalization and the override map}

Left to itself, Claude Code canonicalizes a configured id --- it would send
$gwCanon~g~m$, the wrong thing when that differs from $m$ (a bare
\texttt{claude-opus-4-8} routes first-party, not to the selected provider). The
configuration corrects this with an override map so the verbatim $m$ reaches the
wire. $wireOf$ is the \emph{correct} target: always the id itself.

\begin{axdef}
  wireOf : ModelId \fun ModelId
\where
  \forall m : ModelId @ wireOf~m = m
\end{axdef}

Tier resolution: a tier/small request is served by a tier-appropriate id in the
selection's provider family --- catalogued, entitled, same provider, generally a
different model. $tierId$ models it on a gateway.

\begin{axdef}
  tierId : Gateway \cross ModelId \pfun ModelId
\where
  \forall g : Gateway; m : ModelId @ \\
  \t1 ( ( g, m ) \in \dom tierId \iff m \in gwServes~g ) \land \\
  \t1 ( m \in gwServes~g \implies \\
  \t2   ( tierId~(g, m) \in gwServes~g \land \\
  \t2     gwCatalog~g~(tierId~(g, m)) = gwCatalog~g~m ) )
\end{axdef}

\section{Configuration state}

$sel$ the selection; $selGw$ its gateway; $routed$ whether the gateway is active;
$activeGw$ the gateway in force; $auth$ the auth mode; $wire$ the id per request
kind; $override$ whether the \texttt{modelOverrides} map is present (both the
canonical-key and raw-identity entries Claude Code needs). The selection must be
entitled AND a Claude id --- a non-Claude id is never a valid configured state.

\begin{schema}{Config}
  sel : ModelId \\
  selGw : Gateway \\
  routed : Bool \\
  activeGw : Gateway \\
  auth : AuthMode \\
  wire : ReqKind \pfun ModelId \\
  override : Bool
\where
  sel \in validOn~selGw \\
  sel \in claudeIds
\end{schema}

\section{Initial state}

Never configured: direct, no managed wire, no override.

\begin{schema}{Init}
  Config'
\where
  routed' = no \\
  auth' = vkauth \\
  activeGw' = selGw' \\
  wire' = \emptyset \\
  override' = no
\end{schema}

\section{Well-formedness}

When routed, $wire$ carries exactly Claude Code's three request kinds.

\begin{schema}{WellFormed}
  Config
\where
  routed = yes \implies \dom wire = \{ main, tier, small \}
\end{schema}

\section{The core guarantee}

\textbf{Routing correctness.} Routed, the main request carries the selection
verbatim, the tier and small requests carry a same-family tier id, and the
override is present so canonicalization does not corrupt the main id. The wire
target for kind $k$ under selection $m$ on gateway $g$ is written inline as
\[ \IF k = main \THEN m \ELSE tierId~(g, m), \]
inlined rather than a function over the $ReqKind \cross ModelId$ product so the
model checker enumerates only the three finite kinds.

\begin{schema}{RoutingCorrect}
  WellFormed
\where
  routed = yes \implies \\
  \t1 ( ( \forall k : \dom wire @ \\
  \t2     wire~k = ( \IF k = main \THEN sel \ELSE tierId~(activeGw, sel) ) ) \land \\
  \t1   override = yes )
\end{schema}

\textbf{Provider fidelity.} Every request --- main or side --- reaches the
selection's provider on the active gateway. Entitlement makes that provider
reachable, excluding the catalogued-but-\texttt{AccessDenied} failure.

\begin{schema}{ProviderFidelity}
  RoutingCorrect
\where
  routed = yes \implies \\
  \t1 ( \forall k : \dom wire @ \\
  \t2   provOn~(activeGw, wire~k) = provOn~(activeGw, sel) )
\end{schema}

\textbf{The forbidden failure.} The historical bug: the main request sent the
canonicalized base id (no override), or a side-request fell back to a built-in
whose provider is not the selection's.

\begin{schema}{SilentMisroute}
  WellFormed
\where
  routed = yes \\
  ( override = no \lor \\
  \t1 ( \exists k : \dom wire @ \\
  \t2   provOn~(activeGw, wire~k) \neq provOn~(activeGw, sel) ) )
\end{schema}

\section{Shared effect: applying a gateway configuration}

Route through gateway $g$ with Claude selection $m$: rebuild $wire$ (main $= m$,
tier/small $= tierId~(g,m)$) and set the override.

\begin{schema}{ApplyGateway}
  \Delta Config \\
  g? : Gateway \\
  m? : ModelId
\where
  sel' = m? \\
  selGw' = g? \\
  activeGw' = g? \\
  routed' = yes \\
  auth' = auth \\
  \dom wire' = \{ main, tier, small \} \\
  wire'~main = m? \\
  wire'~tier = tierId~(g?, m?) \\
  wire'~small = tierId~(g?, m?) \\
  override' = yes
\end{schema}

\section{Operation: \texttt{use}}

Selecting a model. Input is a gateway, a model id, and a $force$ flag.
Partitions: (a) valid --- a Claude id, entitled, gateway usable; (b) forced ---
$force? = yes$ selects a Claude id that is catalogued but not entitled, gateway
usable; (c) refused-nonclaude --- the id is not a Claude id (the loud
\texttt{unrecognized\_model} outcome), refused even under force, because Claude
Code cannot run it at all; (d) refused-unknown --- a Claude id not catalogued and
not forced; (e) refused-gateway --- gateway lacks the Anthropic protocol.

\begin{schema}{UseOk}
  ApplyGateway \\
  force? : Bool
\where
  m? \in claudeIds \\
  m? \in validOn~g? \\
  g? \in usableGateways \\
  force? = no
\end{schema}

\begin{schema}{UseForced}
  ApplyGateway \\
  force? : Bool
\where
  force? = yes \\
  m? \in claudeIds \\
  m? \in \dom (gwCatalog~g?) \\
  m? \notin gwServes~g? \\
  g? \in usableGateways
\end{schema}

The loud refusal: a non-Claude id. No state change --- Claude Code cannot run it,
and \texttt{--force} does not help. Distinct code.

\begin{schema}{UseRefusedNonClaude}
  \Xi Config \\
  g? : Gateway \\
  m? : ModelId \\
  force? : Bool \\
  err! : \num
\where
  m? \notin claudeIds \\
  err! = 1
\end{schema}

\begin{schema}{UseRefusedUnknown}
  \Xi Config \\
  g? : Gateway \\
  m? : ModelId \\
  force? : Bool \\
  err! : \num
\where
  m? \in claudeIds \\
  m? \notin \dom (gwCatalog~g?) \\
  force? = no \\
  err! = 2
\end{schema}

\begin{schema}{UseRefusedGateway}
  \Xi Config \\
  g? : Gateway \\
  m? : ModelId \\
  force? : Bool \\
  err! : \num
\where
  m? \in claudeIds \\
  g? \notin usableGateways \\
  err! = 3
\end{schema}

\begin{zed}
  Use \defs (UseOk \lor UseForced) \lor \\
  \t1 (UseRefusedNonClaude \lor \\
  \t2 (UseRefusedUnknown \lor UseRefusedGateway))
\end{zed}

\section{Operation: \texttt{auth}}

Set the auth mode, preserving the current route. $subauth$ (the subscription
credential) is always legitimate here: this tool is Claude Code, Anthropic's own
client, and $usableGateways$ already requires the Anthropic protocol. The only
refusal is a gateway that does not accept the requested mode.

\begin{schema}{AuthOk}
  \Delta Config \\
  newauth? : AuthMode
\where
  sel' = sel \\
  selGw' = selGw \\
  activeGw' = activeGw \\
  routed' = routed \\
  auth' = newauth? \\
  wire' = wire \\
  override' = override
\end{schema}

\begin{zed}
  Auth \defs AuthOk
\end{zed}

\section{Operations: gateway-on, \texttt{direct}, \texttt{toggle}}

\begin{schema}{GatewayOn}
  \Delta Config
\where
  selGw \in usableGateways \\
  sel' = sel \\
  selGw' = selGw \\
  activeGw' = selGw \\
  routed' = yes \\
  auth' = auth \\
  \dom wire' = \{ main, tier, small \} \\
  wire'~main = sel \\
  wire'~tier = tierId~(selGw, sel) \\
  wire'~small = tierId~(selGw, sel) \\
  override' = yes
\end{schema}

\begin{schema}{Direct}
  \Delta Config
\where
  routed' = no \\
  sel' = sel \\
  selGw' = selGw \\
  activeGw' = activeGw \\
  auth' = auth \\
  wire' = \emptyset \\
  override' = no
\end{schema}

\begin{schema}{Toggle}
  \Delta Config
\where
  sel' = sel \\
  selGw' = selGw \\
  ( routed = no \implies \\
  \t1 ( routed' = yes \land activeGw' = selGw \land auth' = auth \land \\
  \t2   selGw \in usableGateways \land \\
  \t2   \dom wire' = \{ main, tier, small \} \land \\
  \t2   wire'~main = sel \land \\
  \t2   wire'~tier = tierId~(selGw, sel) \land \\
  \t2   wire'~small = tierId~(selGw, sel) \land \\
  \t2   override' = yes ) ) \\
  ( routed = yes \implies \\
  \t1 ( routed' = no \land activeGw' = activeGw \land auth' = auth \land \\
  \t2   wire' = \emptyset \land override' = no ) )
\end{schema}

\section{Read-only operations: \texttt{status}, \texttt{list}, \texttt{verify}}

$Verify$ is the probe: $ok! = 1$ exactly when routed, well-formed (all three
request kinds present, so a missing tier/small slot cannot pass vacuously),
every request reaches the selection's provider, and the override is present.

\begin{schema}{Status}
  \Xi Config \\
  repSel! : ModelId \\
  repGw! : Gateway \\
  repRouted! : Bool \\
  repAuth! : AuthMode
\where
  repSel! = sel \\
  repGw! = activeGw \\
  repRouted! = routed \\
  repAuth! = auth
\end{schema}

\begin{schema}{List}
  \Xi Config \\
  cat! : \power ModelId
\where
  cat! = validOn~activeGw \cap claudeIds
\end{schema}

\begin{schema}{VerifyPass}
  \Xi Config \\
  WellFormed \\
  ok! : \num
\where
  routed = yes \\
  ( \forall k : \dom wire @ \\
  \t1 provOn~(activeGw, wire~k) = provOn~(activeGw, sel) ) \\
  override = yes \\
  ok! = 1
\end{schema}

\begin{schema}{VerifyFail}
  \Xi Config \\
  ok! : \num
\where
  ( routed = no \lor \\
  \t1 ( \lnot ( \dom wire = \{ main, tier, small \} ) ) \lor \\
  \t1 ( \exists k : \dom wire @ \\
  \t2   provOn~(activeGw, wire~k) \neq provOn~(activeGw, sel) ) \lor \\
  \t1 override = no ) \\
  ok! = 0
\end{schema}

\begin{zed}
  Verify \defs VerifyPass \lor VerifyFail
\end{zed}

\section{The system invariant}

Routed, the configuration is well-formed, every request carries its correct wire
target and reaches the selection's provider, and the override is present. Stated
as a machine-checkable theorem.

\begin{schema}{SystemOK}
  WellFormed
\where
  routed = yes \implies \\
  \t1 ( ( \forall k : \dom wire @ \\
  \t2     wire~k = ( \IF k = main \THEN sel \ELSE tierId~(activeGw, sel) ) ) \land \\
  \t1   ( \forall k : \dom wire @ \\
  \t2     provOn~(activeGw, wire~k) = provOn~(activeGw, sel) ) \land \\
  \t1   override = yes )
\end{schema}

\begin{theorem}{ApplyGivesSystemOK}
  \forall Config'; g? : Gateway; m? : ModelId @ \\
  \t1 ( \exists Config @ ApplyGateway ) \implies \\
  \t2 ( \theta Config' \in \{ SystemOK @ \theta Config \} )
\end{theorem}

\section{On-disk configuration}

Claude Code's config lives in two files: the managed block in
$DIR$\texttt{/.envrc.local} (environment) and \texttt{.claude/settings.local.json}
(\texttt{modelOverrides}). A slot holds a concrete id or $inherit$. A main slot
\emph{resolves to the selection} ($inherit$ or $fixed~sel$); a tier/small slot
\emph{resolves within the family} ($inherit$ or a $fixed$ entitled id in the
selection's provider family).

\begin{zed}
  Slot ::= inherit | fixed \ldata ModelId \rdata
\end{zed}

\begin{axdef}
  resolves : ModelId \fun \power Slot
\where
  \forall m : ModelId @ \\
  \t1 resolves~m = \{ inherit \} \cup \{ s : Slot | s = fixed~m \}
\end{axdef}

\begin{axdef}
  resolvesFamily : Gateway \cross ModelId \pfun \power Slot
\where
  \forall g : Gateway; m : ModelId @ \\
  \t1 ( ( g, m ) \in \dom resolvesFamily \iff m \in gwServes~g ) \land \\
  \t1 ( m \in gwServes~g \implies \\
  \t2   resolvesFamily~(g, m) = \{ inherit \} \cup \\
  \t3   \{ s : Slot; i : gwServes~g | s = fixed~i \land \\
  \t4     gwCatalog~g~i = gwCatalog~g~m @ s \} )
\end{axdef}

Slots: $ccMain$ (\texttt{ANTHROPIC\_MODEL}, the selection verbatim, no $[1m]$
marker --- the marker breaks the override-map key match); the four tier aliases
$ccOpus$/$ccSonnet$/$ccHaiku$/$ccFable$
(\texttt{ANTHROPIC\_DEFAULT\_\{OPUS,SONNET,HAIKU,FABLE\}\_MODEL}); $ccSmall$
(\texttt{ANTHROPIC\_SMALL\_FAST\_MODEL}); and $ccOverride$ (the
\texttt{modelOverrides} map carrying the canonical-key and raw-identity entries).

\begin{schema}{ClaudeCodeConfig}
  ccMain, ccOpus, ccSonnet, ccHaiku, ccFable, ccSmall : Slot \\
  ccOverride : Bool
\end{schema}

Claude Code is correct when routed iff the main slot resolves to the selection,
the override is present, and every tier/small slot resolves within the family.

\begin{schema}{ClaudeCodeOK}
  ClaudeCodeConfig \\
  sel : ModelId \\
  selGw : Gateway
\where
  sel \in gwServes~selGw \\
  ccMain \in resolves~sel \\
  ccOpus \in resolvesFamily~(selGw, sel) \land ccSonnet \in resolvesFamily~(selGw, sel) \\
  ccHaiku \in resolvesFamily~(selGw, sel) \land ccFable \in resolvesFamily~(selGw, sel) \\
  ccSmall \in resolvesFamily~(selGw, sel) \\
  ccOverride = yes
\end{schema}

\section{Why the hard parts are contained}

Everything in this spec that is not in \texttt{model-provider-fleet} ---
$CanonId$, $gwCanon$, $wireOf$, $tierId$, $resolvesFamily$, the override, the
tier/small request kinds, the $subauth$ mode, the \texttt{unrecognized\_model}
refusal --- exists solely because Claude Code canonicalizes, aliases tiers, and
is Anthropic's own subscription client. Confining it to a single-harness tool
means a change to any OpenAI-protocol harness can never regress this path, and a
change here can never regress the fleet. The complexity is irreducible, but it is
now bounded, named, and model-checked rather than tangled through a tool that
also has to serve three harnesses that need none of it.

\end{document}
